\documentclass[../../main.tex]{subfiles}

\begin{document}

\section{Stone-Čech 紧致化}\label{section:ch05-stone-cech}

我们已经研究过拓扑空间 $X$ 紧致化的一种方法, 那便是单点紧致化(第 29 节). 从某种意义上说, 这是 $X$ 的极小紧致化, 而 Stone-Čech 紧致化则是 $X$ 的极大紧致化. 这里介绍的这种紧致化是由 M. Stone 和 E. Čech 于 1937 年各自独立完成的. 它在近代分析中有着许多应用, 但那些都已超出了本书的范围.

让我们先来回顾以下定义:

\begin{definition}\label{definition:ch05-compactification}
空间 $X$ 的一个紧致化(compactification) $Y$ 是一个包含 $X$ 的紧致的 Hausdorff 空间, 以 $X$ 为其子空间, 并且使得 $\overline{X}=Y$. $X$ 的两个紧致化 $Y_1$ 和 $Y_2$ 是等价的(equivalent), 若存在一个同胚 $h\colon Y_1\to Y_2$, 使得对于每一个 $x\in X$ 有 $h(x)=x$.
\end{definition}

如果空间 $X$ 有一个紧致化 $Y$, 那么 $X$ 必须是完全正则的, 因为它是完全正则空间 $Y$ 的子空间. 反之, 若 $X$ 是完全正则空间, 则 $X$ 至少有一个紧致化. 这是因为对于某一个 $J$, $X$ 能被嵌入到紧致的 Hausdorff 空间 $[0,1]^J$ 中, 并且, 正如下面的引理所说, 任何一个这样的嵌入将会给出 $X$ 的一个紧致化:

\begin{lemma}\label{lemma:ch05-induced-compactification}
设 $X$ 是一个空间, $h\colon X\to Z$ 是从 $X$ 到紧致的 Hausdorff 空间 $Z$ 的一个嵌入. 则存在 $X$ 的一个相应的紧致化 $Y$ 具有以下性质: 存在一个嵌入 $H\colon Y\to Z$ 使得 $H$ 在 $X$ 上的限制等于 $h$. 在不区别等价的两个空间的意义下, 紧致化 $Y$ 是唯一确定的.
\end{lemma}

这个空间 $Y$ 称为由嵌入 $h$ 所诱导的紧致化.

\begin{proof}
给定 $h$, 令 $X_0$ 为 $Z$ 的子空间 $h(X)$, $Y_0$ 为它在 $Z$ 中的闭包. 于是 $Y_0$ 是一个紧致的 Hausdorff 空间, 并且 $\overline{X_0}=Y_0$. 因此 $Y_0$ 就是 $X_0$ 的一个紧致化.

现在我们来构造一个包含 $X$ 的空间 $Y$, 使得空间偶对 $(X,Y)$ 同胚于空间偶对 $(X_0,Y_0)$. 选择与 $X$ 无交的一个集合 $A$, 以及 $A$ 与 $Y_0-X_0$ 之间的一个一一对应 $k\colon A\to Y_0-X_0$. 定义 $Y=X\cup A$, 并且定义一一对应 $H\colon Y\to Y_0$ 为
\begin{align*}
H(x)=h(x),&\quad \text{对于 }x\in X,\\
H(a)=k(a),&\quad \text{对于 }a\in A.
\end{align*}
赋予 $Y$ 一个拓扑, 使得 $U$ 是 $Y$ 的一个开集当且仅当 $H(U)$ 是 $Y_0$ 的一个开集. 映射 $H$ 当然是一个同胚; 而且空间 $X$ 是 $Y$ 的一个子空间, 因为当 $H$ 限制在 $Y$ 的子集 $X$ 上时, 等于同胚 $h$. 通过扩大 $H$ 的值域, 我们便得到了所求的从 $Y$ 到 $Z$ 的一个嵌入.

现在假设 $Y_i$ 是 $X$ 的一个紧致化, 嵌入 $H_i\colon Y_i\to Z$ 是 $h$ 的扩张, $i=1,2$. 于是 $H_i$ 将 $X$ 映射到 $h(X)=X_0$ 上. 因为 $H_i$ 是连续的, 所以它将 $Y_i$ 映射到 $\overline{X_0}$ 中; 又因为 $H_i(Y_i)$ 包含着 $X_0$ 并且还是闭的(因为是紧致的), 所以 $H_i(Y_i)$ 包含着 $\overline{X_0}$. 因此 $H_i(Y_i)=\overline{X_0}$, 并且 $H_2^{-1}\circ H_1$ 定义了从 $Y_1$ 到 $Y_2$ 的一个同胚, 并且在 $X$ 上等于恒等映射.
\end{proof}

一般地说, 对于一个给定的空间 $X$, 可以有多种不同的紧致化方法. 例如对于开区间 $X=(0,1)$, 可有以下一些紧致化:

\begin{examples}
在 $\mathbb{R}^2$ 中取单位圆周 $S^1$, 令映射 $h\colon(0,1)\to S^1$ 为
\begin{align*}
h(t)=(\cos 2\pi t)\times(\sin 2\pi t).
\end{align*}
嵌入 $h$ 诱导的紧致化等价于 $X$ 的单点紧致化.
\end{examples}

\begin{examples}
设 $Y$ 为空间 $[0,1]$. 则 $Y$ 是 $X$ 的一个紧致化, 它是由``在 $(0,1)$ 的两端各加上一点''而得到的.
\end{examples}

\begin{examples}
考虑 $\mathbb{R}^2$ 中单位方形 $[-1,1]^2$, 令映射 $h\colon(0,1)\to[-1,1]^2$ 为
\begin{align*}
h(x)=x\times\sin(1/x).
\end{align*}
空间 $Y_0=\overline{h(X)}$ 为拓扑学家的正弦曲线(见第 24 节例 7). 嵌入 $h$ 所产生的 $(0,1)$ 的紧致化与前面两个完全不同, 它是由在 $(0,1)$ 的右端加上一点, 在左端加上一条线段而得到的.
\end{examples}

紧致化研究中的一个基本问题是:

\begin{quote}
如果 $Y$ 是 $X$ 的一个紧致化, 那么在什么条件下, 定义在 $X$ 上的一个连续实值函数 $f$ 可以连续地扩张到 $Y$ 上?
\end{quote}

函数 $f$ 如果能够扩张, 则它必是有界的, 因为它的扩张将紧致空间 $Y$ 映入 $\mathbb{R}$ 中, 因而是有界的. 但是, 有界性一般说来并不充分, 考虑下面的例子:

\begin{examples}
设 $X=(0,1)$, 考虑例 1 中给出的 $X$ 的单点紧致化. 有界连续函数 $f\colon(0,1)\to\mathbb{R}$ 可以扩张到这个紧致化上的充分必要条件是极限
\begin{align*}
\lim_{x\to 0+}f(x)\quad\text{和}\quad\lim_{x\to 1-}f(x)
\end{align*}
存在并且相等.
\end{examples}

对于例 2 中给出的 $X$ 的``两点紧致化''函数 $f$ 可以扩张的充分必要条件是上述两个极限都存在.

例 3 中的紧致化, 对于一类更广泛的函数都存在扩张. 容易看到, 如果上面两个极限都存在, 则 $f$ 可以扩张. 但是, 函数 $f(x)=\sin(1/x)$ 也可以扩张到这个紧致化上: 令 $H$ 是从 $Y$ 到 $\mathbb{R}^2$ 中的嵌入, 在子空间 $X$ 上它等于 $h$, 则复合映射
\begin{align*}
Y\xrightarrow{H}\mathbb{R}\times\mathbb{R}\xrightarrow{\pi_2}\mathbb{R}
\end{align*}
就是所要求的 $f$ 的扩张. 因为如果 $x\in X$, 则 $H(x)=h(x)=x\times\sin(1/x)$, 于是 $\pi_2(H(x))=\sin(1/x)$.

关于最后这个紧致化, 有一点非常重要, 为了得到它, 我们选择一个嵌入
\begin{align*}
h\colon(0,1)\longrightarrow\mathbb{R}^2,
\end{align*}
其分量函数是 $x$ 及 $\sin(1/x)$. 这时我们发现, 函数 $x$ 和 $\sin(1/x)$ 都能扩张到 $X$ 的紧致化上. 这对于我们很有启发, 如果在 $(0,1)$ 上定义了连续有界函数构成的族, 我们可用它们作为从 $(0,1)$ 到 $\mathbb{R}^J$ 中的一个嵌入(对于某一个 $J$), 从而得到一个紧致化, 使上述族中每一个函数都可以扩张.

这一思想是 Stone-Čech 紧致化中的主要思想, 陈述如下:

\begin{theorem}\label{theorem:ch05-stone-cech-existence}
设 $X$ 是完全正则空间. 则存在 $X$ 的一个紧致化 $Y$ 满足条件: 对于每一个有界连续函数 $f\colon X\to\mathbb{R}$ 都可以唯一地扩张为从 $Y$ 到 $\mathbb{R}$ 的一个连续函数.
\end{theorem}

\begin{proof}
设 $\{f_\alpha\}_{\alpha\in J}$ 是 $X$ 上所有有界连续实值函数所构成的族, 以 $J$ 为其指标集. 对于每一个 $\alpha\in J$, 选取 $\mathbb{R}$ 中一个包含 $f_\alpha(X)$ 的闭区间 $I_\alpha$. 为了确定起见, 选取
\begin{align*}
I_\alpha=[\inf f_\alpha(X),\sup f_\alpha(X)].
\end{align*}
然后定义 $h\colon X\to\prod_{\alpha\in J}I_\alpha$ 为
\begin{align*}
h(x)=(f_\alpha(x))_{\alpha\in J}.
\end{align*}

根据\cref{theorem:ch05-tychonoff}, $\prod I_\alpha$ 是紧致的. 由于 $X$ 是完全正则的, 函数族 $\{f_\alpha\}$ 分离 $X$ 中的点和闭集. 因此根据定理 34.2, $h$ 是一个嵌入.

设 $Y$ 是由嵌入 $h$ 所诱导的紧致化. 则存在一个嵌入 $H\colon Y\to\prod I_\alpha$, 它在 $Y$ 的子空间 $X$ 上的限制就是 $h$. 给定 $X$ 上的一个连续有界实值函数 $f$, 我们指出它可以扩充到 $Y$ 上. 因为函数 $f$ 属于族 $\{f_\alpha\}_{\alpha\in J}$, 所以存在指标 $\beta$, 使得 $f$ 等于 $f_\beta$. 设 $\pi_\beta\colon\prod I_\alpha\to I_\beta$ 为投射. 则连续映射 $\pi_\beta\circ H\colon Y\to I_\beta$ 就是所要求的 $f$ 的扩张. 这是因为, 如果 $x\in X$, 则有
\begin{align*}
\pi_\beta(H(x))=\pi_\beta(h(x))=\pi_\beta\bigl((f_\alpha(x))_{\alpha\in J}\bigr)=f_\beta(x).
\end{align*}
扩张的唯一性是以下引理的一个推论.
\end{proof}

\begin{lemma}\label{lemma:ch05-extension-uniqueness}
设 $A\subset X$, $f\colon A\to Z$ 是从 $A$ 到 Hausdorff 空间 $Z$ 中的一个连续映射. 则最多存在 $f$ 的一个连续扩张 $g\colon\overline{A}\to Z$.
\end{lemma}

\begin{proof}
这个引理曾在第 18 节中作为习题出现过. 我们在这里给出证明. 假定 $g,g'\colon\overline{A}\to X$ 是 $f$ 的两个不同的扩张, 选择 $x$ 使得 $g(x)\neq g'(x)$. 令 $U$ 和 $U'$ 分别为 $g(x)$ 和 $g'(x)$ 的无交的邻域, 选取 $x$ 的邻域 $V$, 使得 $g(V)\subset U$ 并且 $g'(V)\subset U'$. 于是 $V$ 与 $A$ 交于某一点 $y$, 从而 $g(y)\in U$, $g'(y)\in U'$. 但是, 由于 $y\in A$, 所以 $g(y)=f(y)$ 及 $g'(y)=f(y)$. 这就与 $U$ 及 $U'$ 无交矛盾.
\end{proof}

\begin{theorem}\label{theorem:ch05-stone-cech-universal}
设 $X$ 是一个完全正则空间, $Y$ 是 $X$ 的一个紧致化, 它满足\cref{theorem:ch05-stone-cech-existence}所说的扩张条件. 给定从 $X$ 到紧致 Hausdorff 空间 $C$ 中的任意一个连续映射 $f\colon X\to C$, 则 $f$ 可以唯一扩张为一个连续映射 $g\colon Y\to C$.
\end{theorem}

\begin{proof}
注意, 由于 $C$ 是完全正则的, 所以对于某一个 $J$, 它可以嵌入到 $[0,1]^J$ 中. 于是, 我们总可以假定 $C\subset[0,1]^J$. 映射 $f$ 的每一个分支函数 $f_\alpha$ 是 $X$ 上的连续实值有界函数. 根据假设, $f_\alpha$ 可以扩张为从 $Y$ 到 $\mathbb{R}$ 的一个连续映射 $g_\alpha$. 定义 $g\colon Y\to\mathbb{R}^J$ 为 $g(y)=(g_\alpha(y))_{\alpha\in J}$, 我们说映射 $g$ 是连续的, 因为 $\mathbb{R}^J$ 有积拓扑. 实际上 $g$ 确实将 $Y$ 映到 $\mathbb{R}^J$ 的子空间 $C$ 中. 因为根据 $g$ 的连续性, 有
\begin{align*}
g(Y)=g(\overline{X})\subset\overline{g(X)}=\overline{f(X)}\subset\overline{C}=C.
\end{align*}
因此, $g$ 就是所求的 $f$ 的扩张.
\end{proof}

\begin{theorem}\label{theorem:ch05-stone-cech-equivalence}
设 $X$ 是一个完全正则空间. 若 $Y_1$ 和 $Y_2$ 是 $X$ 的两个紧致化, 具有\cref{theorem:ch05-stone-cech-existence}的扩张性质, 则 $Y_1$ 和 $Y_2$ 等价.
\end{theorem}

\par
\begin{figure}[H]
\centering
\includegraphics[width=0.5\textwidth]{fig38-comm.png}
\caption{}\label{figure:ch05-stone-cech-diagram}
\end{figure}

\begin{proof}
考虑内射 $j_2\colon X\to Y_2$, 它是 $X$ 到紧致的 Hausdorff 空间 $Y_2$ 的一个连续映射. 因为 $Y_1$ 具有扩张性质, 由前述定理, 我们可以把 $j_2$ 扩张为一个连续映射 $f_2\colon Y_1\to Y_2$. 类似地, 可以把内射 $j_1\colon X\to Y_1$ 扩张为一个连续映射 $f_1\colon Y_2\to Y_1$(因为 $Y_2$ 具有扩张性质, 而且 $Y_1$ 是一个紧致的 Hausdorff 空间).


复合函数 $f_1\circ f_2\colon Y_1\to Y_1$ 具有以下性质: 对于每一个 $x\in X$, $f_1(f_2(x))=x$. 因此 $f_1\circ f_2$ 是恒等映射 $i_X\colon X\to X$ 的一个连续扩张, 但是 $Y_1$ 上的恒等映射也是 $i_X$ 的一个连续扩张, 按照扩张的唯一性(\cref{lemma:ch05-extension-uniqueness}), $f_1\circ f_2$ 必然等于 $Y_1$ 上的恒等映射. 类似地, $f_2\circ f_1$ 必然等于 $Y_2$ 上的恒等映射. 因此, $f_1$ 与 $f_2$ 是同胚.
\end{proof}

\begin{definition}\label{definition:ch05-stone-cech-beta}
对于每一个完全正则空间 $X$, 取 $X$ 的满足\cref{theorem:ch05-stone-cech-existence}中扩张条件的一个紧致化. 记 $X$ 的这一个紧致化为 $\beta(X)$, 并且称之为 $X$ 的 Stone-Čech 紧致化(Stone-Čech compactification). 它的一个重要性质就是: 从 $X$ 映到一个紧致的 Hausdorff 空间的任意一个连续映射 $f\colon X\to C$ 有唯一的一个连续映射 $g\colon\beta(X)\to C$ 为它的扩张.
\end{definition}

\subsection*{习题}

\begin{enumerate}
\item 证明例 4 中的论断.
\item 证明: 用 $g(x)=\cos(1/x)$ 定义的有界连续函数 $g\colon(0,1)\to\mathbb{R}$ 不能扩张到例 3 中的紧致化. 定义一个嵌入 $h\colon(0,1)\to[0,1]^3$, 使得函数 $x$, $\sin(1/x)$ 及 $\cos(1/x)$ 都能扩张到由 $h$ 所诱导的紧致化上去.
\item 在什么条件下, 可度量化空间有一个可度量化的紧致化?
\item 设 $Y$ 是 $X$ 的任意一个紧致化, $\beta(X)$ 是 Stone-Čech 紧致化. 证明存在一个满的连续闭映射 $g\colon\beta(X)\to Y$, 使得在 $X$ 上它等于恒等映射.

[这个习题清楚地说明 $\beta(X)$ 是 $X$ 的``极大''紧致化. 于是, $X$ 的每一个紧致化都等价于 $\beta(X)$ 的一个商空间.]
\item (a) 证明: 定义在 $S_\Omega$ 上的每一个连续实值函数是``终于常数''的. [提示: 首先证明, 对于每一个 $\varepsilon$, 存在 $S_\Omega$ 的一个元素 $\alpha$ 使得对于所有的 $\beta>\alpha$, 有 $|f(\beta)-f(\alpha)|<\varepsilon$. 然后, 对于 $n\in\mathbb{Z}_+$ 令 $\varepsilon=\tfrac1n$, 再考虑相应点 $\alpha_n$.]
\begin{enumerate}[(b)]
\item 证明: $S_\Omega$ 的单点紧致化和 Stone-Čech 紧致化是等价的.
\item 证明 $S_\Omega$ 的每一个紧致化等价于单点紧致化.
\end{enumerate}
\item 设 $X$ 是一个完全正则空间. 证明: $X$ 是连通的当且仅当 $\beta(X)$ 是连通的. [提示: 如果 $X=A\cup B$ 为 $X$ 的一个分割, 则对于 $x\in A$, 令 $f(x)=0$; 对于 $x\in B$, 令 $f(x)=1$.]
\item 设 $X$ 是一个离散空间, 考虑空间 $\beta(X)$.
\begin{enumerate}[(a)]
\item 证明: 若 $A\subset X$, 则 $\overline{A}$ 与 $\overline{X-A}$ 是无交的, 其中闭包是在 $\beta(X)$ 中求取的.
\item 证明: 若 $U$ 是 $\beta(X)$ 的一个开集, 则 $\overline{U}$ 是 $\beta(X)$ 的一个开集.
\item 证明: $\beta(X)$ 是全不连通的.
\end{enumerate}
\item 证明: $\beta(\mathbb{Z}_+)$ 的基数不小于 $I^I$ 的基数, 其中 $I=[0,1]$. [提示: $I^I$ 有可数稠密子集.]
\item (a) 设 $X$ 是正规的, $y$ 是 $\beta(X)-X$ 的一个点. 证明 $y$ 不是 $X$ 中点的一个序列的极限.
\begin{enumerate}[(b)]
\item 证明: 若 $X$ 是完全正则的但不是紧致的, 则 $\beta(X)$ 不是可度量化的.
\end{enumerate}
\item 对于每一个完全正则的空间 $X$ 与它的 Stone-Čech 紧致化 $\beta(X)$, 我们已经建立了一个对应 $X\to\beta(X)$. 现在对于完全正则空间之间的每一个连续映射 $f\colon X\to Y$, 指派映射 $i\circ f$ 的唯一连续扩张 $\beta(f)\colon\beta(X)\to\beta(Y)$ 与其对应, 其中 $i\colon Y\to\beta(Y)$ 是内射. 验证以下结论:
\begin{enumerate}[(i)]
\item 若 $1_X\colon X\to X$ 是恒等映射, 则 $\beta(1_X)$ 是 $\beta(X)$ 上的恒等映射.
\item 若 $f\colon X\to Y$ 及 $g\colon Y\to Z$, 则 $\beta(g\circ f)=\beta(g)\circ\beta(f)$.
\end{enumerate}
这些性质说明我们所构造的这个对应就是所谓的函子(functor), 它是从完全正则空间和这类空间上连续映射的``范畴'', 到紧致的 Hausdorff 空间和其上连续映射的``范畴''的一个函子. 在本书的第二部分中你将还会看到这些性质; 它们是代数学和代数拓扑学的基础.
\end{enumerate}

\end{document}